
\section{Géométrie des scénarios en boule de Kullback--Leibler}
\label{sec:stress}

Le papier principal retient l'usage le plus simple de cette géométrie : tarifer
en nats un scénario publié, ce qui pour un pur déplacement de moyenne revient à
la moitié du carré de la distance de Mahalanobis. Cette section donne la
construction complète dont cet usage est un cas particulier, avec les
démonstrations que la v1 se contentait d'énoncer, et un résultat nouveau sur le
relèvement d'un scénario de portefeuille à la dimension pleine.

\subsection{Divergence entre deux gaussiennes}

\begin{proposition}[Forme fermée]
\label{prop:klgauss}
Pour $p_1=\mathcal N(\mu_1,\Sigma_1)$ et $p_0=\mathcal N(\mu_0,\Sigma_0)$ sur
$\mathbb R^N$, avec $\delta=\mu_1-\mu_0$,
\begin{equation}
\label{eq:klgauss}
\KL{p_1}{p_0}=\tfrac12\Big[\underbrace{\delta^\top\Sigma_0^{-1}\delta}_{\text{déplacement}}
+\underbrace{\operatorname{tr}(\Sigma_0^{-1}\Sigma_1)-N-\log\det(\Sigma_0^{-1}\Sigma_1)}_{\text{déformation}}\Big].
\end{equation}
\end{proposition}

\begin{proof}
$\KL{p_1}{p_0}=\mathbb E_{p_1}[\log p_1-\log p_0]$. Les deux logarithmes sont
quadratiques ; l'astuce de la trace de l'annexe~\ref{app:gauss} donne
$\mathbb E_{p_1}[(X-\mu_1)^\top\Sigma_1^{-1}(X-\mu_1)]=N$ et
$\mathbb E_{p_1}[(X-\mu_0)^\top\Sigma_0^{-1}(X-\mu_0)]
=\operatorname{tr}(\Sigma_0^{-1}\Sigma_1)+\delta^\top\Sigma_0^{-1}\delta$. On
regroupe avec les termes en $\log\det$.
\end{proof}

Le terme de déplacement est la moitié du carré de la distance de Mahalanobis :
c'est l'unité dans laquelle le papier principal audite les scénarios
publiés. Le terme de déformation est positif et nul si et seulement si
$\Sigma_1=\Sigma_0$ ; en écrivant $\nu_k$ les valeurs propres de
$\Sigma_0^{-1}\Sigma_1$, il vaut $\tfrac12\sum_k(\nu_k-1-\log\nu_k)$, une somme de
termes positifs, un par direction propre.

\subsection{La boule projetée et le taux de change}

\begin{definition}[Scénario en boule d'entropie]
Étant donné la loi conditionnelle d'aujourd'hui $p_0$ et une sévérité $\eta>0$,
le scénario de sévérité $\eta$ pour une fonctionnelle de risque $\rho$ est
\begin{equation}
\label{eq:ball}
p^\star_\eta=\arg\max_{p\,:\,\KL{p}{p_0}\le\eta}\rho(p).
\end{equation}
\end{definition}

Pour un portefeuille linéaire $w^\top r$ de moyenne $m_0$ et d'écart type $s_0$
sous $p_0$, posons $u=(m-m_0)/s_0$ et $v=s/s_0$. Par la
proposition~\ref{prop:klgauss} en dimension un,
\begin{equation}
\label{eq:g}
g(u,v)=\tfrac12\left(u^2+v^2-1-2\log v\right).
\end{equation}

\begin{proposition}[Solution de la boule projetée]
\label{prop:ballsol}
Pour $\rho=-m+c\,s$ avec $c\ge0$ (VaR et expected shortfall gaussiennes sont de
cette forme), le maximum de $\rho$ sous $g(u,v)\le\eta$ est atteint sur la
frontière $g=\eta$ en un point vérifiant
\begin{equation*}
u=-t,\qquad v-v^{-1}=c\,t,\qquad t>0 .
\end{equation*}
Pour $c=0$ : $v=1$ et $u=-\sqrt{2\eta}$. \emph{Un scénario purement
directionnel de $\eta$~nats est un mouvement de $\sqrt{2\eta}$~sigmas} ; un
nat achète $\sqrt2\approx1{,}41$~sigma, et $k$~sigmas coûtent $k^2/2$~nats.
\end{proposition}

\begin{proof}
$g$ est convexe sur $\mathbb R\times(0,\infty)$ (somme de $u^2/2$ et de
$(v^2-1-2\log v)/2$, de dérivée seconde $1+v^{-2}>0$) et l'objectif
$s_0(-u+cv)$ est linéaire : les conditions de Karush--Kuhn--Tucker sont
suffisantes. L'objectif n'a pas de maximum intérieur, donc la contrainte est
saturée. Le lagrangien donne $-1=\lambda u$ et $c=\lambda(v-v^{-1})$ avec
$\lambda>0$ ; en posant $t=1/\lambda$ on obtient le système annoncé, et $t$ est
fixé par $g(-t,v(t))=\eta$. Pour $c=0$, $v-v^{-1}=0$ impose $v=1$ et
$g=u^2/2=\eta$.
\end{proof}

\subsection{Du portefeuille à la dimension pleine}

La boule~\eqref{eq:ball} se résout sur la loi du portefeuille, mais un test de
résistance doit dire comment chaque actif bouge. Deux faits règlent la
question.

\begin{proposition}[Traitement des données]
\label{prop:dpi}
Pour toute loi $p$ sur $\mathbb R^N$, et $p^w$, $p_0^w$ les lois de $w^\top r$
sous $p$ et $p_0$,
$\KL{p^w}{p_0^w}\le\KL{p}{p_0}$.
\end{proposition}

C'est l'inégalité de traitement des données \citep{cover2006} appliquée à
l'application $r\mapsto w^\top r$. Le prix \emph{ambiant} d'un scénario (en
dimension pleine) est donc toujours au moins son prix \emph{vu du portefeuille}
(en projection). Le tableau des bundles classiques de la v1 le vérifie ligne à
ligne ($9{,}97$ contre $7{,}24$~nats pour le bundle standard).

\begin{proposition}[Relèvement minimal]
\label{prop:lift}
Soit $p_0=\mathcal N(\mu_0,\Sigma_0)$, $m_0=w^\top\mu_0$, $s_0^2=w^\top\Sigma_0w$,
et une cible $(m,s)$ avec $s>0$. Parmi les gaussiennes
$\mathcal N(\mu_1,\Sigma_1)$ telles que $w^\top\mu_1=m$ et
$w^\top\Sigma_1w=s^2$, la divergence $\KL{p_1}{p_0}$ est minimale pour
\begin{equation}
\label{eq:lift}
\mu_1=\mu_0+u\,s_0\,\beta,\qquad
\Sigma_1=\Sigma_0+(v^2-1)\,s_0^2\,\beta\beta^\top,
\qquad\beta=\frac{\Sigma_0w}{s_0^2},
\end{equation}
et ce minimum vaut exactement $g(u,v)$, le prix vu du portefeuille.
\end{proposition}

\begin{proof}
Par la proposition~\ref{prop:dpi}, tout candidat coûte au moins le prix de sa
projection, qui est $g(u,v)$ puisque la projection d'une gaussienne est
gaussienne de moyenne $m$ et d'écart type $s$. Il suffit donc de vérifier que
\eqref{eq:lift} atteint cette borne. On a $w^\top\beta=1$, donc
$w^\top\mu_1=m_0+us_0=m$ et $w^\top\Sigma_1w=s_0^2+(v^2-1)s_0^2=s^2$. Dans
\eqref{eq:klgauss} : $\delta^\top\Sigma_0^{-1}\delta=u^2s_0^2\,\beta^\top\Sigma_0^{-1}\beta=u^2$ ;
$\operatorname{tr}(\Sigma_0^{-1}\Sigma_1)=N+(v^2-1)s_0^2\,\beta^\top\Sigma_0^{-1}\beta=N+v^2-1$ ;
et par le lemme du déterminant matriciel,
$\det(\Sigma_0^{-1}\Sigma_1)=1+(v^2-1)s_0^2\,\beta^\top\Sigma_0^{-1}\beta=v^2$.
D'où $\KL{p_1}{p_0}=\tfrac12(u^2+v^2-1-2\log v)=g(u,v)$. La matrice $\Sigma_1$ est
définie positive pour tout $v>0$.
\end{proof}

\paragraph{Lecture actif par actif.} Le vecteur $\beta$ est celui des bêtas des
actifs sur le portefeuille, $\beta_i=\operatorname{Cov}(r_i,w^\top r)/s_0^2$.
Écrivons $r=\beta f+e$ avec $f=w^\top r$ et $e$ non corrélé à $f$ sous $p_0$.
Le scénario~\eqref{eq:lift} multiplie l'écart type du facteur portefeuille $f$
par $v$, déplace sa moyenne de $u$~sigmas, et laisse la loi de $e$ inchangée.
Chaque actif voit sa variance augmenter de $(v^2-1)s_0^2\beta_i^2$ et sa
moyenne bouger de $us_0\beta_i$ : les actifs à fort bêta prennent le choc, et
les corrélations entre eux montent. C'est un scénario à un facteur, et il est
le moins informatif de tous ceux qui atteignent la sévérité visée.

\begin{remark}[Correction de la v1]
La version arXiv~v1 identifiait ce minimum par ses conditions du premier ordre,
$\Sigma_1^{-1}=\Sigma_0^{-1}-\lambda ww^\top$, et le jugeait \og non
interprétable actif par actif \fg{}. Par la formule de Sherman--Morrison, cette
perturbation de rang un de la précision \emph{est}~\eqref{eq:lift}, une
perturbation de rang un de la covariance le long de $\Sigma_0w$, et elle
s'interprète comme ci-dessus. La v1 ne l'utilisant nulle part, aucun chiffre
n'en dépend.
\end{remark}

\begin{corollary}[Coût de la dilatation uniforme]
\label{cor:uniform}
La dilatation uniforme $\Sigma_1=v^2\Sigma_0$, choix par défaut de la v1, a une
composante de déformation égale à $N\cdot\tfrac12(v^2-1-2\log v)$, soit $N$~fois
celle du relèvement minimal.
\end{corollary}

\begin{proof}
Toutes les valeurs propres de $\Sigma_0^{-1}\Sigma_1$ valent $v^2$ ; on applique
la forme spectrale du terme de déformation.
\end{proof}

Pour le scénario décennal du panel principal ($v=1{,}63$, $N=20$), la
composante de variance vaut $0{,}34$~nat au minimum et $20\times0{,}34\approx6{,}8$
sous la dilatation uniforme, ce qui est le chiffre du tableau de composition de
la v1. Le corollaire explique le constat empirique de la v1 selon lequel les
compositions calibrées sur H1 ($2{,}88$~nats) et raffinées par H8 ($0{,}71$)
coûtent \emph{moins} que le défaut : le défaut déforme également toutes les
directions propres, alors qu'une composition qui concentre le choc sur le mode
commun se rapproche du scénario à un facteur~\eqref{eq:lift}, dont le coût total
est $0{,}41$~nat avec le déplacement de moyenne. Que la dilatation uniforme
coûte $N$~fois le minimum en variance est un résultat général ; que les
compositions calibrées se placent entre les deux reste un constat sur ces
panels, mais il devient attendu plutôt que surprenant.

\subsection{Composer une sévérité donnée}

\begin{proposition}[Une sévérité, plusieurs compositions]
\label{prop:composition}
Si $w^\top\Sigma_aw=w^\top\Sigma_bw=s^2$, alors pour tout $\alpha\in[0,1]$,
$\Sigma_\alpha=(1-\alpha)\Sigma_a+\alpha\Sigma_b$ vérifie
$w^\top\Sigma_\alpha w=s^2$.
\end{proposition}

\begin{proof}
$w^\top\Sigma w$ est linéaire en $\Sigma$.
\end{proof}

La sévérité du portefeuille ne contraint donc pas la composition. Le prix
ambiant, lui, est convexe en $\Sigma_1$ (le terme $-\log\det$ est convexe et le
reste est linéaire), donc le long d'un mélange il est majoré par le mélange des
prix, ce qui garantit qu'une composition intermédiaire ne coûte jamais plus que
la plus chère de ses deux extrémités.

\paragraph{Un plafond physique pour la corrélation.} À volatilités marginales
fixées, $w^\top\Sigma w$ est maximal sur les matrices de corrélation valides pour
$\rho_{ij}=1$, où il vaut $(\sum_iw_i\sigma_i)^2$. Sur le panel principal ce
plafond ne vaut qu'environ $1{,}5$~fois la volatilité d'aujourd'hui : au-delà de
l'échelon quinquennal, aucune hausse de corrélation seule ne suffit, et une part
de hausse des volatilités est inévitable.

\subsection{Calibrer le rayon, et la limite de cette calibration}

Le rayon se lit sur le flux d'information de la projection,
$I^{(h)}_t=\KL{p^w_t}{p^w_{t-h}}$, dont les quantiles (sur blocs non
chevauchants de 21~jours) convertissent une période de retour en rayon. Lire le
rayon sur le flux à $N$ dimensions le surestime d'un ordre de grandeur, par la
proposition~\ref{prop:dpi}. Calibrée ainsi, l'échelle de la v1 donne un
multiplicateur de volatilité décennal de $1{,}63$, quand le quantile décennal de
la hausse de volatilité réalisée sur un mois vaut $3{,}2$. L'échelle est une
règle cohérente et reproductible ; ses périodes de retour ne sont pas validées.
C'est l'une des raisons pour lesquelles le papier principal garde l'audit
(tarifer un scénario donné) et abandonne la génération (construire un scénario
d'une période de retour donnée).

\subsection{Le test de résistance inversé}

Demander le scénario le moins coûteux qui produit une perte $L$ donnée inverse
la boule. Dans le cadre gaussien projeté, la
proposition~\ref{prop:ballsol} se lit à l'envers. Dans le cadre non
paramétrique, la solution est l'inclinaison exponentielle de l'échantillon
\citep{csiszar1975}, dont la taille d'échantillon effective s'effondre quand le
rayon croît : à $9$~nats, elle vaut $1{,}0$ sur le panel principal, c'est-à-dire
que la loi inclinée n'est plus que la pire journée observée. Toute
implémentation devrait rapporter cette taille effective et basculer vers la
boule gaussienne dès qu'elle tombe sous quelques dizaines.
